% Generated draft: requires mathematical and source review.
\documentclass[11pt]{article}
\usepackage[margin=25mm]{geometry}
\usepackage{fontspec}
\setmainfont{Latin Modern Roman}
\usepackage{amsmath,amssymb,mathtools,unicode-math}
\setmathfont{Latin Modern Math}
\usepackage{enumitem,needspace,xurl,hyperref}
\setlength{\parindent}{0pt}
\setlength{\parskip}{5pt}
\setlength{\emergencystretch}{3em}
\newcommand{\N}{\mathbb{N}}
\newcommand{\Z}{\mathbb{Z}}
\newcommand{\Q}{\mathbb{Q}}
\newcommand{\R}{\mathbb{R}}
\newcommand{\C}{\mathbb{C}}
\newcommand{\F}{\mathbb{F}}
\newcommand{\A}{\mathbb{A}}
\newcommand{\PP}{\mathbb{P}}
\newcommand{\E}{\mathbb{E}}
\newcommand{\eps}{\varepsilon}
\newcommand{\dd}{\,\mathrm{d}}
\begin{document}
\section*{UnsolvedMath ID: 30000580}
\textbf{Source problem code:} catalogue entry
\subsection{Definitions and mathematical statement}
Under the usual complex-algebraic convention, let \(X\) be smooth and complete, let \(D\subset X\) be a reduced simple-normal-crossings divisor, and let a connected algebraic group \(G\) act preserving \(D\). The pair is log homogeneous under \(G\) when the infinitesimal-action map
\[
\mathcal O_X\otimes_{\C}\operatorname{Lie}(G)\longrightarrow T_X(-\log D)
\]
is surjective, where the logarithmic tangent sheaf consists of vector fields tangent to the boundary. Its open homogeneous part is the intended \(G/H=X\setminus D\). The source does not reproduce the base-field, completeness, or group-class conventions, so these standard definitions need source matching; in particular \(G\) is not silently required to be reductive or linear.

The first question is to classify, without a dimension bound, homogeneous spaces \(G/H\) admitting a log homogeneous equivariant completion \((X,D)\). A completion here contains the homogeneous space as a dense open orbit and extends its \(G\)-action, with the boundary and logarithmic-generation property above. The exact isomorphism and group-action conventions of the requested classification must be fixed.

For a pair in the intended class, let \(\tau_X\) be its Tits map. Write its Stein factorization as
\[
X\xrightarrow{\sigma}S(X)\xrightarrow{\nu}\tau_X(X),
\qquad\tau_X=\nu\circ\sigma,
\]
with \(\sigma\) having connected fibers and \(\nu\) finite. The source does not specify its Tits-map construction or target; a chosen ordinary or logarithmic isotropy construction cannot be imported silently.

The second question is whether \(S(X)\) always has only quotient singularities. Over \(\C\), a standard local analytic meaning is that each point has a neighborhood germ isomorphic to the quotient of a smooth germ by a finite group. An algebraic or \emph{\'etale}-local convention, if intended instead, should be specified. This is a condition on \(S(X)\), not a claim that every completion \(X\) is itself a quotient, and no weaker rational-singularity criterion is substituted.

The supplied literature assessment checked 2026-08-21 reports classification and quotient-singularity criteria for important spherical and horospherical classes, but not the arbitrary-dimensional classification requested in the report. Those special classes do not by their labels cover the general log-homogeneous setting.
\subsection{Short English statement}
Classify homogeneous spaces admitting smooth equivariant completions with globally generated logarithmic tangent sheaf, and ask whether the Stein factor of the intended Tits map always has quotient singularities. The two tasks and the general group scope are kept separate.
\subsection{Sources}
\begin{itemize}
\item Local UnsolvedMath registry, permanent ID 30000580, catalogue entry, \emph{Equivariant Completions of Log-Homogeneous Spaces}.
\item \emph{Komplexe Analysis} (2007), Oberwolfach Reports, pp. 2399--2446, DOI \url{https://doi.org/10.4171/owr/2006/40}.
\item Supplied reference ``quotient singularities of horospherical varieties'': \url{https://doi.org/10.1007/s00031-026-09959-x}.
\item Supplied literature triage checked 2026-08-21; no exact Tits construction or complete special-case theorem hypotheses are reproduced.
\end{itemize}
\end{document}
